% Método Ideal:
% - Descrição Analítica
% - Simulação de amostragem
% - Filtro de reconstrução
% - Simulação de reconstrução

% Método Natural:
% - Descrição Analítica
% - Simulação de amostragem
% - Filtro de reconstrução
% - Simulação de reconstrução

% Método Topo-Plano:
% - Descrição Analítica
% - Simulação de amostragem
% - Filtro de reconstrução
% - Simulação de reconstrução

\section{Desenvolvimento analítico} %talvez mudar o nome dessa seção 

\begin{frame}[allowframebreaks]{Parâmetros fixos}
    \small

    \textbf{Sinal original e espectro:}
    \begin{align*}
        x(t) &= \cos(2\pi \cdot 2t) + 0,5\cos(2\pi \cdot 18t) \\
        &\Updownarrow \\
        X(f) &= \frac{1}{2}[\delta(f-2) + \delta(f+2)]
        + \frac{0,5}{2}[\delta(f-18) + \delta(f+18)]
    \end{align*}

    \textbf{Filtro Anti-aliasing ($H_{AA}$):}
    \begin{align*}
        X_f(f) &= X(f) \cdot H_{AA}(f)
        \longrightarrow
        H_{AA}(f) =
        \begin{cases}
            1, & |f| \le 8\,\text{Hz}\\
            0, & |f| > 8\,\text{Hz}
        \end{cases}
        \quad \leftarrow f_c \\[2pt]
        X_f(f) &= \frac{1}{2}[\delta(f-2) + \delta(f+2)]
        \longleftrightarrow
        x_f(t) = \cos(2\pi \cdot 2t)
    \end{align*}

    \textbf{Parâmetros de Amostragem:}
    \[
        f_{max} = 2\,\text{Hz}
        \implies
        f_s > 2f_{max} = 4\,\text{Hz}
    \]

    \[
        f_s = 5\,\text{Hz}
        \implies
        T_s = \frac{1}{5}\,\text{s}
        \implies
        \text{banda de guarda} = 5 - 4 = 1\,\text{Hz}
    \]

    A partir dessa preparação, utiliza-se a equação filtrada em banda para realizar as amostragens: \begin{align*} x_f(t) = \cos(2\pi \cdot 2t) \end{align*}
    
\end{frame}

\begin{frame}[fragile]{Amostragem ideal}
    Para a amostragem ideal, multiplica-se o sinal por um trem de impulsos:
    $$X_s(f) = f_s \cdot \sum_{k=-\infty}^{\infty} X_f(f - k f_s) = 5 \cdot \sum_{k=-\infty}^{\infty} X_f(f - 5k)$$
    
    \textbf{Análise das componentes:}
    $$
    \left.\begin{matrix}
    k=0: \pm 2\text{ Hz} \\
    k=1: 3 \text{ e } 7\text{ Hz} \\
    k=-1: -7 \text{ e } -3\text{ Hz} \\
    k=2: 8 \text{ e } 12\text{ Hz}
    \end{matrix}\right\} \text{componentes em } \pm 2, \pm 3, \pm 7, \pm 8, \pm 12, \dots
    $$
\end{frame}
    

